\documentclass{beamer}
\usepackage[utf8]{inputenc}

\usetheme{Madrid}
\usecolortheme{default}
\usepackage{amsmath,amssymb,amsfonts,amsthm}
\usepackage{mathtools}
\usepackage{txfonts}
\usepackage{tkz-euclide}
\usepackage{listings}
\usepackage{adjustbox}
\usepackage{array}
\usepackage{gensymb}
\usepackage{tabularx}
\usepackage{gvv}
\usepackage{lmodern}
\usepackage{circuitikz}
\usepackage{tikz}
\lstset{literate={·}{{$\cdot$}}1 {λ}{{$\lambda$}}1 {→}{{$\to$}}1}
\usepackage{graphicx}

\setbeamertemplate{page number in head/foot}[totalframenumber]

\usepackage{tcolorbox}
\tcbuselibrary{minted,breakable,xparse,skins}



\definecolor{bg}{gray}{0.95}
\DeclareTCBListing{mintedbox}{O{}m!O{}}{%
  breakable=true,
  listing engine=minted,
  listing only,
  minted language=#2,
  minted style=default,
  minted options={%
    linenos,
    gobble=0,
    breaklines=true,
    breakafter=,,
    fontsize=\small,
    numbersep=8pt,
    #1},
  boxsep=0pt,
  left skip=0pt,
  right skip=0pt,
  left=25pt,
  right=0pt,
  top=3pt,
  bottom=3pt,
  arc=5pt,
  leftrule=0pt,
  rightrule=0pt,
  bottomrule=2pt,
  toprule=2pt,
  colback=bg,
  colframe=orange!70,
  enhanced,
  overlay={%
    \begin{tcbclipinterior}
    \fill[orange!20!white] (frame.south west) rectangle ([xshift=20pt]frame.north west);
    \end{tcbclipinterior}},
  #3,
}
\lstset{
    language=C,
    basicstyle=\ttfamily\small,
    keywordstyle=\color{blue},
    stringstyle=\color{orange},
    commentstyle=\color{green!60!black},
    numbers=left,
    numberstyle=\tiny\color{gray},
    breaklines=true,
    showstringspaces=false,
}
%------------------------------------------------------------
%This block of code defines the information to appear in the
%Title page
\title %optional
{7.4.33}
\date{September 8,2025}
%\subtitle{A short story}

\author % (optional)
{Harsha-EE25BTECH11026}



\begin{document}


\frame{\titlepage}


\begin{frame}{Question}
A circle C of radius 1 unit is inscribed in an equilateral triangle PQR. The points of contact of C with sides PQ, QR, RP are D, E, F respectively.The line PQ is given by the equation $\sqrt{3}x+y-6=0$ and the point $\vec{D}$ is $\brak{\frac{3\sqrt{3}}{2},\frac{3}{2}}$. Further, it is giventhat the origin and the centre of C are on same side of line PQ. The equation of circle C is
\begin{enumerate}
    \item $\brak{x-2\sqrt{3}}^2+\brak{y-1}^2=1$
    \item $\brak{x-2\sqrt{3}}^2+\brak{y+\frac{1}{2}}^2=1$
    \item $\brak{x-\sqrt{3}}^2+\brak{y-1}^2=1$
    \item $\brak{x-\sqrt{3}}^2+\brak{y+1}^2=1$
\end{enumerate}
\end{frame}

\begin{frame}{Theoretical Solution}
According to the question,
\begin{align}
    \text{Equation of tangent PQ :\,}\vec{n}^{\top}\vec{x}=c
\end{align}
where $\vec{n}=\myvec{\sqrt{3}&&1}^{\top}$ and $c=6$
\begin{align}
    \text{Point of tangency}\brak{\vec{D}}:\myvec{\frac{3\sqrt{3}}{2}\\\frac{3}{2}}
\end{align}
\begin{align}
    radius(r)=1
\end{align}
\end{frame}

\begin{frame}{Theoretical Solution}
From the given data,
\begin{align}
    \vec{D}-\vec{u}=\lambda\vec{n}\;\text{, for some scalar $\lambda$}
\end{align}
\begin{align}
    \implies \vec{u}=\vec{D}-\lambda\vec{n}
\end{align}
Also,
\begin{align}
    \frac{|\vec{n}^{\top}\vec{O}-c|}{\|\vec{n}\|}=r
\end{align}
\begin{align}
    |\vec{n}^{\top}\vec{u}-c|=r\|\vec{n}\|
\end{align}
\begin{align}
    \vec{n}^{\top}\vec{u}=c \pm r\|\vec{n}\|
\end{align}
\end{frame}

\begin{frame}{Theoretical Solution}
To decide the sign , we need to use the fact that the origin and centre of circle are on the same side of the line PQ.
\begin{align}
    \therefore \brak{\vec{n}^{\top}\vec{u}-c}\brak{\vec{n}^{\top}\myvec{0\\0}-c}>0
\end{align}
\begin{align}
    \implies \vec{n}^{\top}\vec{u}<c
\end{align}
\begin{align}
    \therefore \vec{n}^{\top}\vec{u}=c-r\|\vec{n}\|
\end{align}
Substituting value of $\vec{u}$,
\begin{align}
\vec{n}^{\top}\brak{\vec{D}-\lambda\vec{n}}=c-r\|\vec{n}\|
\end{align}
\begin{align}
    \implies \lambda=\frac{\vec{n}^{\top}\vec{D}+r\|n\|-c}{\vec{n}^{\top}\vec{n}}
\end{align}
\begin{align}
    \vec{u}=\vec{D}-\frac{\vec{n}^{\top}\vec{D}+r\|n\|-c}{\vec{n}^{\top}\vec{n}}\;\vec{n}
\end{align}
\end{frame}

\begin{frame}{Theoretical Solution}
Substituting the values,
\begin{align}
    \therefore \vec{u}=\myvec{\frac{3\sqrt{3}}{2}\\\frac{3}{2}}
-\frac{1}{2}\myvec{\sqrt{3}\\1}=\myvec{\sqrt{3}\\1}
\end{align}
\begin{align}
    \therefore \text{Required equation of circle :\,}\|\vec{x}\|^2-2\myvec{\sqrt{3}&&1}\vec{x}+3=0
\end{align}
\end{frame}

\begin{frame}[fragile]
    \frametitle{C Code -Finding the equation of circle}

    \begin{lstlisting}[language=C]
#include <stdio.h>
#include <math.h>

void circle_equation(double a, double b, double c, double xd, double yd, double r) {
    double norm = sqrt(a*a + b*b);
    double nx = a / norm;
    double ny = b / norm;
    
    double c1x = xd + r * nx;
    double c1y = yd + r * ny;
    double c2x = xd - r * nx;
    double c2y = yd - r * ny;

    double origin_side = a*0 + b*0 + c;
    double c1_side = a*c1x + b*c1y + c;
    double c2_side = a*c2x + b*c2y + c;
    double h, k;


    \end{lstlisting}
\end{frame}

\begin{frame}[fragile]
    \frametitle{C Code -Finding the equation of circle}

    \begin{lstlisting}[language=C]
    if ((origin_side >= 0 && c1_side >= 0) || (origin_side <= 0 && c1_side <= 0)) {
        h = c1x; k = c1y;
    } else {
        h = c2x; k = c2y;
    }

    // f = h^2 + k^2 - r^2
    double f = h*h + k*k - r*r;

    printf("Equation of circle: (x - %.2lf)^2 + (y - %.2lf)^2 = %.2lf^2\n", h, k, r);
    
}
    \end{lstlisting}
\end{frame}

\begin{frame}[fragile]
    \frametitle{Python+C code}

    \begin{lstlisting}[language=Python]
import ctypes
import numpy as np
import matplotlib.pyplot as plt
import matplotlib as mp
mp.use("TkAgg")
# Load C library
lib = ctypes.CDLL("./circle_equation.so")
# Define argument types
lib.circle_equation.argtypes = [ctypes.c_double, ctypes.c_double, 
ctypes.c_double,ctypes.c_double, ctypes.c_double, ctypes.c_double]
# Define return type (void)
lib.circle_equation.restype = None

a, b, c = np.sqrt(3), 1, -6
xd, yd = 3*np.sqrt(3)/2, 3/2
r = 1
lib.circle_equation(a, b, c, xd, yd, r)

    \end{lstlisting}
\end{frame}

\begin{frame}[fragile]
    \frametitle{Python+C code}

    \begin{lstlisting}[language=Python]
#plotting
h, k = np.sqrt(3), 1
theta = np.linspace(0, 2*np.pi, 500)
x_circle = h + r*np.cos(theta)
y_circle = k + r*np.sin(theta)
# Line PQ
x_vals = np.linspace(0, 6, 200)
y_vals = -(a*x_vals + c)/b
# Tangency point D
D = np.array([xd, yd])
plt.figure(figsize=(6,6))
plt.plot(x_circle, y_circle, label="Circle C")
plt.plot(x_vals, y_vals, label="Line PQ")
plt.scatter(*D, color="red", label="Tangency D")
plt.scatter(h, k, color="black", label="Center O")
    \end{lstlisting}
\end{frame}

\begin{frame}[fragile]
    \frametitle{Python+C code}
    \begin{lstlisting}[language=Python]
plt.axis("equal")
plt.grid(True, linestyle="--", alpha=0.6)
plt.legend()
plt.title("Circle Plot")
plt.savefig("/home/user/Matrix/Matgeo_assignments/7.4.33/figs/Figure_1")
plt.show()
    \end{lstlisting}   
\end{frame}

\begin{frame}[fragile]
    \frametitle{Python code}
    \begin{lstlisting}[language=Python]
import numpy as np
import matplotlib.pyplot as plt
import sympy as sp
import matplotlib as mp
mp.use("TkAgg")

x, y = sp.symbols('x y', real=True)
sqrt3 = sp.sqrt(3)

O = sp.Matrix([sqrt3, 1])
r = sp.Integer(1)
X=np.array([x,y])

expr = (X.dot(X) - 2*O.dot(X) + (O.dot(O) - r**2)).expand()
sp.pprint(expr)
print("=0")
    \end{lstlisting}   
\end{frame}

\begin{frame}[fragile]
    \frametitle{Python code}
    \begin{lstlisting}[language=Python]
h = float(sqrt3.evalf())
k = 1.0
r_num = 1.0

# --- plotting ---
theta = np.linspace(0, 2*np.pi, 500)
x_circle = h + r_num * np.cos(theta)
y_circle = k + r_num * np.sin(theta)

a, b, c = np.sqrt(3), 1.0, -6.0

x_vals = np.linspace(h - 3, h + 6, 400)
y_vals = -(a*x_vals + c) / b

D = np.array([3*np.sqrt(3)/2, 3/2])
plt.figure(figsize=(7,7))

    \end{lstlisting}   
\end{frame}

\begin{frame}[fragile]
    \frametitle{Python code}
    \begin{lstlisting}[language=Python]
plt.plot(x_circle, y_circle, label="Circle C")
plt.plot(x_vals, y_vals, label=r"Line PQ: $\sqrt{3}x + y - 6 = 0$")
plt.scatter([D[0]], [D[1]], color='red', zorder=5, label="D (tangency)")
plt.scatter([h], [k], color='black', zorder=5, label=r"Center $O=(\sqrt{3},1)$")
plt.gca().set_aspect('equal', 'box')
plt.grid(True, linestyle='--', alpha=0.6)
plt.legend()
plt.title("Equation of Circle C")
plt.xlabel("x")
plt.ylabel("y")
plt.savefig("/home/user/Matrix/Matgeo_assignments/7.4.33/figs/Figure_1.png")
plt.show()

    \end{lstlisting}   
\end{frame}

\begin{frame}{Plot}
    \begin{figure}[H]
    \centering
    \includegraphics[width=0.6\columnwidth]{figs/Figure_1.png}
    \label{fig:1}
\end{figure}
\end{frame}

\end{document}